\documentclass[a4paper, sigconf]{acmart}
\usepackage{multirow,hhline,graphicx,array}
\usepackage{standalone}
\newcolumntype{M}[1]{>{\centering\arraybackslash}m{#1}}

%\title{Black-Box Transformation of the GASLITE Attack}
%\author{Ishay Yemini\\Tel Aviv University\\ \today}
%\date{\today}

\settopmatter{printacmref=false}
\setcopyright{none}
\renewcommand\footnotetextcopyrightpermission[1]{}
\pagestyle{plain}

\begin{document}

\appendix
\twocolumn

\section{Combination Attack}


\subsection{Overview}
%While prior work, such as \texttt{GASLITE} \cite{bentov2024}, demonstrated the ability to manipulate retrieval rankings via gradient-based methods, these "white-box" attacks relied on unrealistic access to the internal architecture of the model.
%To bridge this gap, the Combination Attack is presented as a novel black-box methodology. This approach successfully generates adversarial passages that hijack the retrieval process, without requiring access to any internal model architecture, such as gradients, parameters or internal weights - instead relying only on the tokenizer and embedding outputs.
%
%The main achievement of this methodology is the demonstration that black-box attacks can achieve performance parity with state-of-the-art white-box attacks, like \texttt{GASLITE} \cite{bentov2024}. The Combination Attack consistently achieves perfect top-1 retrieval success rates across a diverse suite of modern open- and closed-source models. % TODO, specific which?

While gradient-based methods like \texttt{GASLITE} \cite{bentov2024} can manipulate retrieval rankings, they require "white-box" access to model internals. 
The Combination Attack bridges this gap as a black-box methodology, hijacking the retrieval process using only tokenizer and embedding outputs. 
Crucially, this approach achieves performance parity with white-box states-of-the-art, consistently reaching perfect top-1 success rates across diverse open- and closed-source models.


\subsection{Methodology}

The Combination Attack is a multi-stage optimizer that appends an adversarial "trigger" suffix to a malicious passage ("info"), making it semantically indistinguishable from a target passage in vector space. It uses a hybrid approach to maximize similarity through three distinct strategies:

\subsubsection*{Phase 0: Hot-Start}

Before iterative optimization, the attack employs a "Hot-Start" initialization to anchor the adversarial passage in the target's semantic domain. 
An auxiliary Large Language Model (LLM) is queried to generate 10-20 tokens contextually relevant to the target query. 
This provides a good initial state, significantly reducing the number of steps required in later phases by starting the search in a high-similarity region of the embedding space.


\subsubsection*{Phase 1: Random Search}

%Starting with the hot-start trigger, the algorithm generates a pool of random candidate tokens from the model's vocabulary.
%For each candidate, the algorithm calculates the similarity between the adversarial passage with the appended token and the target passage. 
%The token with the highest similarity is permanently appended, grounding the adversarial passage in the promising semantic region.

Starting from the hot-start seed, the algorithm generates a pool of random candidate tokens from the model's vocabulary. 
Each candidate is evaluated based on the resulting cosine similarity between the adversarial and target passages. 
The highest-scoring token is permanently appended, grounding the sequence in a promising semantic region.


\subsubsection*{Phase 2: Square-like Attack}

%The final stage refines the trigger using a robust, block-base optimization algorithm adapted from the "Square Attack" \cite{andriushchenko2020} used in computer vision.
%Instead of square patches of pixels, this adaptation operated on contiguous segments of text tokens. 
%The algorithm randomly selects a contiguous segment of tokens within the trigger and generates multiple candidate replacements, then if a candidate block increases the similarity score against the target passage, the change is accepted - otherwise, it's ignored. 
%As the attack progresses, the size of the token segment being perturbed decreases. This allows for board exploration in the early iterations, and precise fine-tuning in later stages, essentially "snaking" around the semantic region to align perfectly with the target passage's embedding. 

The final stage refines the sequence using a block-based optimization adapted from the Square Attack \cite{andriushchenko2020}. 
Instead of pixel patches, this adaptation operates on contiguous token segments. 
The algorithm iteratively mutates random segments; replacements are accepted only if they increase the cosine similarity against the target embedding. 
By decreasing the perturbation window over time, the attack transitions from broad exploration to precise semantic alignment.


%\subsection{Implementation Details}
%
%To ensure reproducibility, we provide the specific configurations and auxiliary tools used in the Combination Attack and RASLITE tests. All tests were run with \texttt{seed = 42}.
%
%
%\subsubsection*{Hot-Start Configuration}
%
%The "Hot-Start" was executed using GPT-5-nano \cite{openai2025gpt5nano} as the auxiliary model. 
%The model was prompted using the following instruction:
%
%\textit{"Return a short sentence, up to 15 words, related to the following passage: [Target Passage]"}
%
%
%\subsubsection*{Hyperparameters - Combination Attack}
%
%\begin{itemize}
%  \item \texttt{random\_num\_pool}: 500
%  \item \texttt{random\_early\_stop\_patience}: 5
%  \item \texttt{square\_p\_init}: 0.5
%  \item \texttt{square\_random\_pool\_per\_pos}: 300
%\end{itemize}
%
%
%\subsubsection*{Hyperparameters - RASLITE}
%
%\begin{itemize}
%  \item \texttt{num\_steps}: 1500
%  \item \texttt{n\_candidates}: 128
%  \item \texttt{n\_flip}: 1
%  \item \texttt{n\_bulk\_flips}: 1
%  \item \texttt{flip\_pos\_method}: "ordered"
%\end{itemize}


\subsection{Experimental Evaluation}

The efficacy of the Combination Attack was evaluated against nine popular open-source embedding-based models and three state-of-the-art closed-source embedding models. The results are averaged over 50 trials; in each trial, a unique random query ("target passage") and a random malicious passage ("info") were chosen.  

\subsubsection*{Open-Source Models}

\mbox{MiniLM} \cite{wang2020}, \mbox{aMPNet} \cite{song2020}, \mbox{Arctic} \cite{merrick2024}, E5 \cite{wang2024}, \mbox{GTR-T5} \cite{ni2021}, \mbox{Contriever} \cite{izacard2022}, \mbox{Contriever-MS} \cite{izacard2022}, \mbox{ANCE} \cite{xiong2020}, \mbox{mMPNet} \cite{song2020}.


\subsubsection*{Closed-Source Models} 

OpenAI's \mbox{text-embedding-3-small} (OAI-Small) \cite{openai2024embeddings}, OpenAI's \mbox{text-embedding-3-large} (OAI-Large) \cite{openai2024embeddings}, and Google's \mbox{gemini-embedding-001} (Gem-Large) \cite{team2023gemini}.  % TODO more models??


\subsubsection*{Datasets} 

MSMARCO \cite{bajaj2018} for the queries and corpus, and ToxiGen \cite{hartvigsen2022} for the malicious "info" passages. 


\subsubsection*{Baselines} 

We tested \textbf{\texttt{info} Only} where the adversarial passage is just the chosen \texttt{info} alone; \textbf{\texttt{stuffing}}, where the query was appended to the end of \texttt{info}; and \textbf{\texttt{RASLITE}}. % TODO write more about RASLITE?


\begin{table}[H]
\caption{Success Rate (appeared@1) Across Embedding Models. \textnormal{This table compares the top-1 retrieval success of our proposed Combination Attack against standard baselines and \texttt{RASLITE}.}}
 \label{table:res}
 \documentclass[a4paper, sigconf]{acmart}
\begin{document}
\begin{tabular}{l|cccc}
\hline
& \multicolumn{4}{c}{appeared@1} \\
\textbf{Model} & \texttt{info} Only & stuffing & \texttt{RASLITE} & \textbf{Combi.} \\ \hline
MiniLM & 0\% & 10\% & \textbf{100\%} & \textbf{100\%} \\ \hline
aMPNet & 0\% & 18\% & \textbf{100\%} & \textbf{100\%} \\ \hline
Arctic & 0\% & 38\% & \textbf{100\%} & \textbf{100\%} \\ \hline
E5 & 0\% & 2\% & \textbf{100\%} & \textbf{100\%} \\ \hline
GTR-T5 & 0\% & 8\% & 76\% & \textbf{100\%} \\ \hline
Contriever & 0\% & 90\% & \textbf{100\%} & \textbf{100\%} \\ \hline
Contriever-MS & 0\% & 10\% & \textbf{100\%} & \textbf{100\%} \\ \hline
ANCE & 0\% & 0\% & \textbf{100\%} & \textbf{100\%} \\ \hline
mMPNet & 0\% & 30\% & \textbf{100\%} & \textbf{100\%} \\ \hline
OAI-Small\footnotemark[1] & 0\% & 0\% & \textbf{100\%} & \textbf{100\%} \\ \hline
OAI-Large\footnotemark[1] & & & & \\ \hline
Gem-Large\footnotemark[1] & & & & \\ \hline
\end{tabular}
\end{document}

\end{table}

The experimental results displayed in Table \ref{table:res} unequivocally demonstrate the superior performance of the Combination Attack. Across all models, the proposed methodology outperforms all baselines.


\begin{table}[H]
\caption{Token Count Comparison. \textnormal{This table details the total tokens required for Combination Attack and \texttt{RASLITE}, with reduction calculated as the ratio of \texttt{RASLITE} tokens to Combi. tokens.}}
 \label{table:tokens}
 \documentclass[a4paper, sigconf]{acmart}
\begin{document}
\begin{tabular}{l|cc|c}
\hline
& \multicolumn{2}{c|}{Tokens} & \\
\textbf{Model} & \texttt{RASLITE} & \textbf{Combi.} & \textbf{Reduction} ($\times$) \\ \hline
MiniLM & 1,218,702 & \textbf{124,343} & 9.80\\ \hline
aMPNet & 3,019,949 & \textbf{250,816} & 12.04\\ \hline
Arctic & 1,360,054 & \textbf{202,426} & 6.72\\ \hline
E5 & 1,892,159 & \textbf{242,962} & 7.79\\ \hline
GTR-T5 & 3,394,535 & \textbf{278,040} & 12.21\\ \hline
Contriever & 701,259 & \textbf{73,204} & 9.58\\ \hline
Contriever-MS & 810,559 & \textbf{83,794} & 9.67\\ \hline
ANCE & 1,607,391 & \textbf{229,577} & 7.00\\ \hline
mMPNet & 1,570,016 & \textbf{143,010} & 10.98\\ \hline
OAI-Small\footnotemark[1] & 2,556,917 & \textbf{714,794} & 3.58\\ \hline
OAI-Large\footnotemark[1] & & \\ \hline
Gem-Large\footnotemark[1] & & \\ \hline
\hline 
\textbf{\textit{Average}} & 1,813,154 & \textbf{234,296} & 7.74\\ \hline
\end{tabular}
\end{document}

\end{table}


Additionally, the Combination Attack manages to perform the optimization at a much lower cost compared to \texttt{RASLITE}, on average (across all models) having a reduction factor of 7.74$\times$.

\footnotetext[1]{For the closed-source models, the best passage was estimated by using an index on an open-source model to retrieve the top-50 passages, which were then re-embedded to find the maximum similarity; this approach was used to avoid the large cost of indexing the entire corpus with closed-source APIs.}

\bibliographystyle{acm} 
\bibliography{final_paper.bib}


\end{document}

